\sezione{9}{Esperimento 3: traceroute}{da pc2 a pc1}

\subsection{Obiettivo e procedura}
Scoprire \textbf{quali router} attraversa un pacchetto da pc2 a pc1 e capire \emph{come} traceroute ricava questa informazione. wireshark è già collegato alla LAN C: nel browser si cattura su \texttt{eth1}. Su pc2:
\begin{terminale}{pc2}
§root@pc2:/#§ traceroute 195.11.14.5 -z 1
traceroute to 195.11.14.5 (195.11.14.5), 30 hops max, 60 byte packets
 1  ¤200.1.1.1¤ (200.1.1.1)  0.882 ms  0.662 ms  0.456 ms      °# eth0 di r2°
 2  ¤100.0.0.9¤ (100.0.0.9)  0.903 ms  0.877 ms  1.218 ms      °# eth1 di r1°
 3  ¤195.11.14.5¤ (195.11.14.5)  0.987 ms  1.354 ms  1.015 ms  °# eth0 di pc1°
\end{terminale}
\begin{esempio}[Come si legge]
\begin{itemize}
  \item \texttt{30 hops max}: prova al massimo fino a TTL 30. \texttt{60 byte packets}: dimensione delle sonde.
  \item Ogni riga è un \textbf{salto}: il numero è il TTL usato, poi l'indirizzo che ha risposto e tre tempi (una per ciascuna delle 3 sonde).
  \item L'indirizzo mostrato per un router è quello dell'interfaccia \textbf{da cui ha inviato la risposta} verso pc2: per r2 è \texttt{200.1.1.1} (eth0, lato C), per r1 è \texttt{100.0.0.9} (eth1, lato B). Non sono le interfacce «verso pc1»!
  \item \texttt{-z 1}: attende 1 secondo tra una sonda e l'altra, così in Wireshark i pacchetti sono ben separati.
\end{itemize}
\end{esempio}

\subsection{Il trucco del TTL}
\begin{intuizione}
Ogni pacchetto IP ha un \textbf{TTL} (\emph{Time To Live}), un contatore che ogni router decrementa di 1. Se arriva a 0, il router butta il pacchetto e avvisa il mittente con un messaggio ICMP \textbf{Time Exceeded} (\emph{time-to-live exceeded in transit}), che porta il proprio indirizzo come sorgente. Serve a impedire che i pacchetti girino all'infinito in caso di cicli di routing.

traceroute sfrutta questo meccanismo «di proposito»: manda sonde con TTL = 1, poi 2, poi 3\dots{} La sonda con TTL = $k$ «muore» esattamente al $k$-esimo router, che si rivela rispondendo.
\end{intuizione}
\begin{center}
\rowcolors{2}{chiaro}{white}
\begin{tabularx}{\linewidth}{c l Y l}
\rowcolor{navy}\thd{TTL sonda} & \thd{Dove si ferma} & \thd{Cosa succede} & \thd{Risposta}\\
1 & r2 & r2 decrementa a 0, scarta la sonda & ICMP Time Exceeded da 200.1.1.1\\
2 & r1 & r2 porta il TTL a 1, r1 a 0 e scarta & ICMP Time Exceeded da 100.0.0.9\\
3 & pc1 & arriva a pc1 con TTL 1: è la destinazione & ICMP Port Unreachable da 195.11.14.5\\
\end{tabularx}
\end{center}
\begin{tecnica}[Perché la destinazione risponde «Port Unreachable»]
Le sonde di traceroute (versione Linux, default) sono datagrammi \textbf{UDP} indirizzati a pc1 su porte alte (a partire da 33434) dove con ogni probabilità nessun programma è in ascolto. Quando una sonda arriva davvero a pc1, il kernel di pc1 risponde con ICMP \textbf{Destination Unreachable / Port Unreachable}. Per traceroute questo è il segnale che la destinazione è stata raggiunta e che può fermarsi.
\end{tecnica}

\subsection{La cattura sulla LAN C}
\begin{center}
\footnotesize\ttfamily
\setlength{\tabcolsep}{4pt}
\begin{tabularx}{\linewidth}{r l l l l >{\raggedright\arraybackslash}X}
\rowcolor{navy}\thd{No.} & \thd{Source} & \thd{Destination} & \thd{Proto} & \thd{TTL} & \thd{Info}\\
\wsrow{1}{200.1.1.7}{195.11.14.5}{UDP}{1}{33434 $\rightarrow$ ... sonda 1/3}{blu!8}
\wsrow{2}{200.1.1.1}{200.1.1.7}{ICMP}{64}{Time-to-live exceeded (in transit)}{rosso!20}
\wsrow{\dots}{}{}{}{}{\textrm{altre 2 sonde con TTL=1 e relative risposte}}{white}
\wsrow{7}{200.1.1.7}{195.11.14.5}{UDP}{2}{sonda con TTL=2}{blu!8}
\wsrow{8}{100.0.0.9}{200.1.1.7}{ICMP}{63}{Time-to-live exceeded (in transit)}{rosso!20}
\wsrow{\dots}{}{}{}{}{\textrm{altre 2 sonde con TTL=2}}{white}
\wsrow{13}{200.1.1.7}{195.11.14.5}{UDP}{3}{sonda con TTL=3}{blu!8}
\wsrow{14}{195.11.14.5}{200.1.1.7}{ICMP}{62}{Destination unreachable (Port unreachable) \textrm{$\leftarrow$ destinazione raggiunta!}}{rosso!12}
\end{tabularx}
\end{center}
\didascalia{Ricostruzione illustrativa (numerazione indicativa). In Wireshark i messaggi ICMP di errore sono evidenziati con righe scure. Il campo \emph{Time to live} della sonda si legge nel dettaglio del pacchetto (header IP): 1, poi 2, poi 3. Durante il dialogo compaiono anche ARP unicast di verifica delle cache.}
\begin{attenzione}[Le righe scure in Wireshark non sono errori di cattura]
Wireshark colora in nero/rosso i pacchetti ICMP di errore (Time Exceeded, Destination Unreachable). Qui sono \emph{attesi}: sono proprio le risposte su cui si basa traceroute.
\end{attenzione}

\subsection{Il quadro completo}
\begin{center}
\begin{tikzpicture}[y=-1cm,x=1cm,
  msg/.style={-{Stealth[length=5pt]},line width=1pt},
  ml/.style={font=\ttfamily\scriptsize,fill=white,inner sep=1pt}]
  \foreach \x/\n/\s in {0/pc1/host,4.6/r1/rtr,9.2/r2/rtr,13.8/pc2/host}{
    \node[\s,minimum size=.85cm,minimum width=1.2cm] at (\x,0) {\n};
    \draw[grigio!60,line width=1pt] (\x,.5) -- (\x,8.0);
  }
  \foreach \x/\l in {2.3/A,6.9/B,11.5/C}{\node[cd,minimum width=.7cm,minimum height=.5cm] at (\x,0) {\l};}
  \draw[-{Stealth},grigio] (-.9,.8) -- node[sloped,above,font=\headfont\scriptsize]{tempo} (-.9,7.8);
  % ttl 1
  \draw[msg,blu] (13.8,1.1) -- node[ml,above]{udp, ttl=1} (9.2,1.1);
  \draw[msg,rosso] (9.2,1.7) -- node[ml,above]{icmp ttl exceeded} (13.8,1.7);
  \node[font=\headfont\bfseries\small,text=viola] at (14.6,1.4) {$\times3$};
  % ttl 2
  \draw[msg,blu] (13.8,3.0) -- node[ml,above]{udp, ttl=2} (9.2,3.0);
  \draw[msg,blu] (9.2,3.0) -- node[ml,above]{udp, ttl=1} (4.6,3.0);
  \draw[msg,rosso] (4.6,3.6) -- node[ml,above]{icmp ttl exceeded} (13.8,3.6);
  \node[font=\headfont\bfseries\small,text=viola] at (14.6,3.3) {$\times3$};
  % ttl 3
  \draw[msg,blu] (13.8,5.0) -- node[ml,above]{udp, ttl=3} (9.2,5.0);
  \draw[msg,blu] (9.2,5.0) -- node[ml,above]{udp, ttl=2} (4.6,5.0);
  \draw[msg,blu] (4.6,5.0) -- node[ml,above]{udp, ttl=1} (0,5.0);
  \draw[msg,rosso] (0,5.6) -- node[ml,above,pos=.16]{icmp port unreachable} (13.8,5.6);
  \node[font=\headfont\bfseries\small,text=viola] at (14.6,5.3) {$\times3$};
  \node[draw=ambra,fill=ambra!12,rounded corners=3pt,font=\headfont\scriptsize,align=center,text width=5.6cm] at (6.9,7.0) {tutte le sonde UDP hanno come destinazione l'indirizzo IP di \textbf{pc1}: è il TTL a decidere dove si fermano};
\end{tikzpicture}
\end{center}
\didascalia{Per ogni valore di TTL partono 3 sonde ($\times3$). Il TTL indicato sulle frecce è quello del pacchetto in ciascun tratto, dopo i decrementi dei router.}

\begin{mlbox}[Varianti utili]
\begin{itemize}
  \item \inl{traceroute -n 195.11.14.5}: niente risoluzione dei nomi (output più rapido).
  \item \inl{traceroute -I 195.11.14.5}: sonde ICMP Echo invece di UDP; la destinazione risponde con Echo Reply.
  \item Un asterisco \texttt{*} al posto del tempo significa «nessuna risposta entro il timeout»: router che non invia ICMP, filtro, o rotta di ritorno mancante.
\end{itemize}
\end{mlbox}

\begin{appunti}{Riepilogo dell'esperimento 3}
traceroute = sonde UDP verso la destinazione con TTL 1, 2, 3\dots{} (3 sonde per TTL). Ogni router che porta il TTL a 0 risponde \textbf{ICMP Time Exceeded} (rivela il suo indirizzo, quello dell'interfaccia da cui risponde). La destinazione risponde \textbf{ICMP Port Unreachable} (porta UDP chiusa): fine. Output pc2$\to$pc1: 200.1.1.1 (r2), 100.0.0.9 (r1), 195.11.14.5 (pc1). \texttt{-z 1}: 1\,s tra le sonde.
\end{appunti}
